\headline{{\textbf{\Large\sffamily Potting Shed Gossip:}}\\
          {\textbf{\large\sffamily Building a Mighty Trunk}}}
\label{2026-01-gossip}
\noindent
There is a common saying among veterans of the club: you can grow branches in a few seasons, but you can't buy time for a trunk.
A thick, tapering base is the hallmark of a specimen tree, providing that elusive sense of \textit{mochikomi}---the quality of dignity that only comes with age.
If your tree looks more like a spindly sapling than a forest giant, it is time to talk about the ``Sacrifice Branch.''

The mechanics of thickening a trunk are rooted in simple biology.
To increase the girth of the base, the tree must process a high volume of energy.
This is where the \textbf{Sacrifice Branch} comes in.
By allowing a low\--growing branch to grow unchecked for several seasons---often reaching several feet in length---you create a ``superhighway'' of sap flow.
As this branch thickens to support its own massive weight, it pulls nutrients through the main trunk, forcing the wood below it to expand rapidly.

However, power comes with a price.
The key is placement.
A sacrifice branch should ideally be located toward the back of the tree or at a point where the resulting large scar can be hidden or carved into \textit{jin}.
If you let it grow too long or too high, you risk ruining the taper or creating an unsightly ``reverse taper'' bulge.

For particularly stubborn specimens, some practitioners utilize \textbf{trunk piercing or drilling} to trigger an emergency healing response.
By strategically drilling small holes through the cambium at the base, you force the tree to produce a localized swell of callus tissue.
As these wounds heal, the overlapping scar tissue merges to add physical mass and a rugged, weathered texture to the bark.
This ``controlled trauma'' requires a tree in peak health and sterile tools to prevent rot, but it serves as a powerful catalyst for a trunk that has hit a growth plateau.

For those working with deciduous species, another potent technique is \textbf{Ground Growing} or using an oversized ``wooden grow box.''
By allowing the roots to escape into the earth, the tree enters a state of unrestricted growth.
In this ``escape'' phase, the trunk diameter can double in a fraction of the time it would take in a ceramic pot.

When the desired thickness is finally achieved, the sacrifice is made: the large branch is removed.
This leaves a raw wound that, over time, heals into a subtle muscularity.
It requires patience to watch your refined tree look like a wild mess for a year or two, but the result is a powerful, convincing base that anchors the entire composition.
Remember, you aren't just growing a plant; you are thickening a legacy.

\vfill
\begin{center}$\cdots$\end{center}

\noindent
\textit{Got a bonsai\--related rumour, quirky story, or fun fact?}\\
\textit{Share it with us for the next edition of \textit{Potting Shed Gossip}!}
\closearticle
\clearpage
